\begin{center}\textbf{Abstract}\end{center}
\vspace{-6pt}
{\small\noindent\itshape
Property-Based Testing (PBT) checks general properties over automatically generated inputs, but evidence on its effectiveness comes from single ecosystems, and mutation tools generate different mutants for the same program. This work compared the effectiveness, measured by mutation score, of equivalent PBT suites in two equivalent implementations of an event-ticketing web API, one in Java/Spring Boot (jqwik, PIT) and one in Node.js/Express (fast-check, StrykerJS). A controlled factorial experiment crossed the ecosystem, the strength of the properties (eight mostly bound-checking properties versus a suite that adds nine exact-oracle properties) and the number of tries per property (10, 100 and 1000), with 30 runs per configuration (360 runs), each tool's default mutation operators and a harmonized operator set; differences were tested with Mann-Whitney and Vargha-Delaney effect sizes. Oracle strength dominated: with 100 tries, the mutation score rose from 50.0\% to 84.1\% in Java and from 52.1\% to 95.0\% in Node.js. More tries helped only the strong suite. With default operators, Node.js scored higher in every configuration, because PIT mutates return values that state-based properties do not observe; with harmonized operators and 1000 tries, both reached 96.9\%. The ecosystem therefore matters less than how properties are written and which operators are counted.

\noindent\textbf{Keywords:} Test oracle. Test suite quality. Controlled experiment. Mutation tools.\par}

\begin{center}\textbf{Resumo}\end{center}
\vspace{-6pt}
{\small\noindent\itshape
Testes Baseados em Propriedades (PBT, do inglês \textup{Property-Based Testing}) verificam propriedades gerais sobre entradas geradas automaticamente, mas a evidência sobre sua efetividade vem de um único ecossistema, e ferramentas de mutação diferentes geram mutantes diferentes para o mesmo programa. Este trabalho comparou a efetividade, medida pelo \textup{mutation score}, de suítes PBT equivalentes em duas implementações equivalentes de uma interface de programação web de venda de ingressos, uma em Java/Spring Boot (jqwik, PIT) e outra em Node.js/Express (fast-check, StrykerJS). Um experimento controlado fatorial cruzou o ecossistema, a força das propriedades (oito propriedades que verificam sobretudo limites contra uma suíte que acrescenta nove com oráculo exato) e o número de tentativas por propriedade (10, 100 e 1000), com 30 execuções por configuração (360 execuções), os operadores de mutação padrão de cada ferramenta e um conjunto harmonizado; as diferenças foram testadas com Mann-Whitney e tamanho de efeito de Vargha-Delaney. A força do oráculo foi o fator dominante: com 100 tentativas, o \textup{mutation score} subiu de 50,0\% para 84,1\% no Java e de 52,1\% para 95,0\% no Node.js. Mais tentativas só ajudaram a suíte forte. Com os operadores padrão, o Node.js ficou acima em todas as configurações, porque o PIT muta valores de retorno que propriedades baseadas em estado não observam; com operadores harmonizados e 1000 tentativas, os dois chegaram a 96,9\%. O ecossistema pesa menos do que a forma de escrever as propriedades e os operadores considerados.

\noindent\textbf{Palavras-chave:} Oráculo de teste. Qualidade de suítes de teste. Experimento controlado. Ferramentas de mutação.\par}
